% Reconstruction de la page 005
% Meyrueis au fil du temps — Paul Loupiac

\sectiondate{Avant Propos}{}

\vspace{0.7em}

Ces quelques notes ont été rassemblées au cours des dernières années, pour mon usage personnel. Il s’agissait d’abord de garder en mémoire des conversations à bâtons rompus, des souvenirs oraux, parfois écrits, glanés dans des articles de revues ou des archives familiales et remontant souvent au début des années 1900.

Je pensais me borner à quelques feuilles destinées à mes enfants et petits enfants. Mais des parents et amis ayant gardé de fortes attaches à Meyrueis m’ont fait part de leur intérêt. De plus la lecture des « Itinéraires Protestants en Languedoc » et particulièrement en Cévennes, documents et articles réunis par Patrick Cabanel, m’ont incité à aller voir plus précisément dans les Archives de la Lozère mais surtout de l’Hérault (où l’on peut consulter le fonds de l’Ancienne Intendance du Languedoc), les pièces originales concernant des arrestations ou interrogatoires de ce petit peuple de Meyrueis et des environs, déjà « rebroussiers » refusant de se plier aux ordres du pouvoir central : à l’époque l’ordre de Louis XIV.

En effet par où commencer ? Je me suis fixé les années 1500-1600, période où, avec le développement de l’imprimerie, commencent à circuler plus facilement les idées nouvelles dans tous les domaines, religieux, scientifiques ou littéraires. C’est, dans les années 1540-1550 que le protestantisme s’est implanté à Meyrueis : des colporteurs apportaient des livres « luthériens », comme on disait en ce temps là, libelles, traités, chansons spirituelles et psaumes… ; les premiers cultes se seraient tenus vers 1552 dans la bibliothèque du château de Roquedols.

Un texte de 1558, relevé dans les papiers de l’Eglise Catholique de Meyrueis indique que la population est « \textit{infestée de la mâle (mauvaise) opinion} ». C’est semble-t-il, fin 1559 que la ville procède à sa Réformation en décrétant l’abolition de la messe puisque la presque totalité de la population était passée à la Réforme. L’Église Réformée s’organise et reçoit son premier pasteur François Thérond, originaire de Sauve. Il doit desservir en plus de Meyrueis, Gatuzières, Lanuéjols, Montjardin et Saint Sauveur des Pourcils. Deux ans plus tard on lui demande de « dresser » l’Eglise Réformée de Marvejols en plein Gévaudan\footnote{Philippe Chambon dans « Itinéraires Protestants en Cévennes » p.70}.

Mais la paix ne dura guère puisque se déclenchent très vite « les guerres de religion » (il y en eut huit !) entre catholiques et protestants. Elles commencent après le massacre de Wassy en Champagne où une soixantaine de protestants (hommes, femmes et enfants) furent massacrés pour la première fois, par les troupes du Duc de Guise, à l’issue d’un culte ayant rassemblé dans une grange de nombreux fidèles.

\begin{quote}
\itshape
« Les Meyrueisiens prirent une part active aux opérations militaires dont la région fut le théâtre en cette fin du XVI\textsuperscript{o} siècle jusqu’à l’Edit de Nantes (13 avril 1598). Après une relative accalmie les troubles reprirent vers 1620 pour ne plus s’arrêter, en ce qui concerne les persécutions, jusqu’à la veille de la révolution, avec l’édit de 1787 rendant aux non catholiques une existence légale »\footnote{Philippe Chambon dans « Itinéraires Protestants en Cévennes » p.70}.
\end{quote}


Mais il n’était pas question de se limiter à cette période du 17\textsuperscript{e}-18\textsuperscript{e} siècle, marquée par l’arrivée de la Réforme, ses progrès rapides suivis par des mesures de restriction des libertés de plus en plus draconiennes à Meyrueis même, aboutissant à La Révocation de l’Edit de Nantes en 1685 avec la destruction du temple et les abjurations forcées de la plupart des meyrueisiens subissant les dragonnades.

Arrive la période des philosophes et des Lumières qui débouchera sur la Révolution, vécue comme une libération par les protestants du sud de la Lozère. Ne créèrent-ils pas à Meyrueis une « Société Populaire » dès 1793, où se font déjà jour, d’ailleurs, les luttes entre Girondins et Montagnards ? Révolution combattue par les Catholiques, dans la haute Lozère et les Gorges du Tarn, surtout à la suite de la Constitution Civile du Clergé. Il y eut bien des excès de part et d’autre. Mais aux morts républicains répondirent de nombreuses arrestations et condamnations de royalistes. La Malène en souffrit particulièrement car nombre d'hommes de ce village des gorges du Tarn s’étaient engagés dans les bandes royalistes entraînées par Charrier ancien député de la Constituante et 37 furent fusillés à Florac. Il y est fait allusion dans les pages qui suivent.

Vous ne trouverez pas grand chose sur le 19\textsuperscript{e} siècle. Il faudrait y revenir…

Enfin nous arrivons à la période contemporaine, celle que nous avons connue ou que nos pères ont connue au cours des années 1900. Elles commencent avec l'Affaire Dreyfus et les décrets de séparation des Eglises et de l’Etat en 1905.

Des anciens nous ont parlé ou même ont laissé par écrit, des souvenirs précieux, relatés dans cette brochure qui n’est que l’ébauche de ce que pourrait être un ouvrage sinon exhaustif, du moins plus complet, mieux documenté et présentant une véritable iconographie. Nous appelons de nos vœux, une véritable et complète histoire de Meyrueis présentant la vie et les activités de « Meyrueis, ma villetto » depuis les origines jusqu’à nos jours.

\vspace{0.35em}

\begin{tcolorbox}[
  colback=white,
  colframe=black,
  boxrule=0.7pt,
  arc=0pt,
  left=0.25cm,
  right=0.25cm,
  top=0.12cm,
  bottom=0.12cm,
  boxsep=0pt
]
\small
Les pages qui suivent ont été rédigées au cours des quatre ou cinq dernières années, au fur et à mesure, du dépouillement de documents consultés aux Archives de Mende, ou de Montpellier. D’autres renseignements ont été recueillis dans plusieurs livres et publications à la Bibliothèque de la Société d’Histoire du Protestantisme Français (SHPF) à Paris ainsi qu’à la Médiathèque de Montpellier. Enfin sont présentés quelques documents inédits.

C’est la raison pour laquelle vous trouverez des répétitions, et des renseignements qui ne se suivent pas toujours de façon chronologique.
\end{tcolorbox}


